%%%%%%%%%%%%%%%%%%%%%%%%%%%%%%%%%%%%%%%%%%%%%%%%%%%%%%%%%%%%%%%%%%%%%%%%%%%%%%%%
%% 04_Projeto_Fisico_Rede.tex
%% Projeto Físico da Rede, Cabeamento Estruturado e Infraestrutura de Racks
%%%%%%%%%%%%%%%%%%%%%%%%%%%%%%%%%%%%%%%%%%%%%%%%%%%%%%%%%%%%%%%%%%%%%%%%%%%%%%%%

\chapter{Projeto Físico da Rede e Cabeamento Estruturado}
\label{cap:projeto_fisico}

\section{Subsistemas de Cabeamento Estruturado}

A instalação física distribui-se entre os prédios CT-09 e CT-12, estruturada em cinco subsistemas funcionais:
\begin{enumerate}
    \item Área de Trabalho (\textit{Work Area});
    \item Cabeamento Horizontal (\textit{Horizontal Cabling});
    \item Salas de Telecomunicações e Racks de Equipamentos;
    \item Backbone de Campus em Fibra Óptica;
    \item Sistema de Aterramento e Alimentação Elétrica Ininterrupta.
\end{enumerate}

A Figura~\ref{fig:planta_fisica} ilustra o layout físico das salas e a distribuição das tomadas de rede.

\incluirdiagrama{Imagens/planta_fisica_laboratorios.png}{Planta baixa executiva e disposição de pontos de telecomunicações nas salas CT-09 e CT-12.}{fig:planta_fisica}

\section{Área de Trabalho e Conectividade}

A área de trabalho compreende a terminação entre as tomadas modulares instaladas no mobiliário técnico e os computadores dos pesquisadores.

\begin{itemize}
    \item \textbf{Módulos de Conexão}: Emprego de conectores fêmea tipo Keystone RJ-45 Categoria 6 da linha Furukawa SohoPlus ou equivalente homologado, instalados em espelhos modulares embutidos nas bancadas de pesquisa;
    \item \textbf{Cabos de Manobra (\textit{Patch Cords})}: Cabos conectorizados de fábrica com botas injetadas de alívio de tensão mecânica, condutores de cobre multifilar 24 AWG e pinagem T568A em ambas as extremidades. O projeto padroniza comprimentos de 1,5 metro para as estações de bancada e 2,5 metros para pontos de piso ou impressoras;
    \item \textbf{Pontos de Rede Sem Fio}: Instalação de caixas de sobrepor em teto com módulos Keystone Cat6 nos pontos 901 (CT-09) e 1201/WF01 (CT-12), conectando os pontos de acesso TP-Link Omada alimentados via PoE+ (Power over Ethernet, IEEE 802.3at).
\end{itemize}

\section{Cabeamento Horizontal e Infraestrutura de Encaminhamento}

O cabeamento horizontal interliga as tomadas da área de trabalho aos painéis de manobra (\textit{patch panels}) instalados nos racks de telecomunicações de cada prédio.

A extensão máxima do cabo de par trançado no enlace permanente (\textit{Permanent Link}) é limitada a 90 metros. A margem complementar de 10 metros é reservada para os cabos de manobra da área de trabalho e do rack, respeitando o canal total de 100 metros.

O meio físico de transmissão emprega cabos de par trançado não blindados (U/UTP) Categoria 6 de 4 pares, com condutores sólidos de cobre eletrolítico nu de 23 AWG, separador interno em formato de cruz para redução de paradiafonia (NEXT) e capa com baixa emissão de fumaça e livre de halogênios (LSZH).

O encaminhamento físico contempla a desativação das canaletas plásticas sobrepostas preexistentes:
\begin{itemize}
    \item \textbf{Tubulações Primárias}: Eletrodutos de PVC rígido roscável embutidos na alvenaria para descidas verticais até as bancadas;
    \item \textbf{Calhas Secundárias}: Eletrocalhas metálicas perfuradas integradas ao mobiliário dos laboratórios, mantendo separação física mínima de 20 centímetros em relação aos circuitos de energia elétrica para prevenir interferências eletromagnéticas.
\end{itemize}

\section{Dimensionamento dos Racks de Telecomunicações}

Os equipamentos ativos e passivos de rede são acomodados em dois gabinetes metálicos padrão 19 polegadas:
\begin{enumerate}
    \item \textbf{Rack Principal (CT-12 Sala 201)}: Gabinete de piso autossustentado fechado com padrão $42\text{U} \times 800\text{mm} \times 1000\text{mm}$, equipado com porta frontal de vidro temperado, fechadura com chave, painéis laterais ventilados removíveis e teto com módulo de ventilação forçada de 4 ventiladores. Abriga os servidores Dell PowerEdge, os roteadores corporativos e firewall Cisco, os switches de acesso e distribuição, e o nobreak senoidal. O roteador do provedor (\texttt{ISP-UFS}), retratado no domínio externo da topologia lógica, localiza-se no centro de dados central do campus da UFS (NTI), não ocupando espaço físico no rack local;
    \item \textbf{Rack Secundário (CT-09 Sala 101)}: Gabinete padrão 19 polegadas de $24\text{U} \times 600\text{mm} \times 670\text{mm}$, abrigando o switch de acesso local e os painéis de manobra dos 49 pontos de dados do Prédio CT-09.
\end{enumerate}

A Figura~\ref{fig:layout_rack} apresenta a disposição vertical dos módulos no rack principal.

\incluirdiagrama{Imagens/layout_rack_telecomunicacoes.png}{Diagrama de ocupação vertical de equipamentos no Rack Principal de 42U.}{fig:layout_rack}

A Tabela~\ref{tab:layout_rack} sintetiza a distribuição das unidades de rack (RU) no gabinete principal.

\begin{table}[htbp]
\centering
\footnotesize
\caption{Ocupação vertical e distribuição de equipamentos no Rack Principal 42U (CT-12).}
\label{tab:layout_rack}
\begin{tabularx}{\textwidth}{|c|l|Y|Y|}
\hline
\rowcolor{tableheaderbeige}
\textbf{Posição (U)} & \textbf{Dispositivo} & \textbf{Modelo de Referência} & \textbf{Função} \\ \hline
41 -- 42 & DIO Óptico & Furukawa 1U 24 Fibras & Terminação do enlace óptico do campus \\ \hline
39 -- 40 & Roteador de Borda & Cisco ISR 4321 & Conexão WAN UFS e NAT de saída \\ \hline
38 & Firewall Físico & Cisco ASA 5506-X 1U & Inspeção de estado (MPF) e segurança perimetral \\ \hline
36 -- 37 & Roteador Interno & Cisco 2911 & Roteamento inter-VLAN e DHCP Relay \\ \hline
34 -- 35 & Switch Central & Cisco Catalyst 2960-24TT & Distribuição e agregação local CT-12 \\ \hline
32 -- 33 & Guia de Cabos & Organizador Horizontal 1U & Manobra dos cabos patch cord \\ \hline
29 -- 31 & Patch Panels Cat6 & Furukawa SohoPlus 24P & Terminação do cabeamento horizontal \\ \hline
27 -- 28 & Switch PoE+ Wi-Fi & TP-Link Omada SG2210P & Alimentação dos Access Points \\ \hline
22 -- 23 & Servidor DHCP & Servidor 1U (Dell R250/R350) & Serviços DHCP dedicados, DNS e NTP \\ \hline
18 -- 21 & Servidor Web & Dell PowerEdge R440 / R450 1U & Hospedagem NextGIS Web e Nginx \\ \hline
14 -- 17 & Servidor de Banco & Dell PowerEdge R540 / R550 2U & PostgreSQL, PostGIS e Storage \\ \hline
01 -- 04 & Sistema de Energia & Nobreak Senoidal Online 3000VA & Suprimento elétrico ininterrupto \\ \hline
-- & Régua PDU Lateral & PDU 12 Tomadas 20A & Distribuição elétrica com disjuntor \\ \hline
\end{tabularx}
\end{table}

\section{Interligação de Campus em Fibra Óptica}

A conexão entre a Sala 101 do CT-09 e a Sala 201 do CT-12 é realizada por um enlace de fibra óptica monomodo (\textit{Single-Mode Fiber} - SMF) de 9/125 $\mu$m. O cabo óptico possui estrutura geleada com proteção mecânica contra umidade e roedores, apropriado para passagem em dutos subterrâneos entre os blocos universitários.

As fibras são conectorizadas com padrão LC-UPC e ancoradas nos Distribuidores Internos Ópticos (DIO) em ambos os racks. A conversão de mídia óptica é executada através de transceivers SFP Gigabit Ethernet (1000BASE-LX) instalados nas portas de uplink dos switches, garantindo tráfego imune a interferências eletromagnéticas.

\section{Identificação e Administração do Cabeamento}

Todos os componentes físicos da rede recebem identificação padronizada com etiquetas em poliéster térmico segundo a sintaxe:
\[
\text{[Prédio]}-\text{[Sala]}-\text{[Rack]}-\text{[Painel]}-\text{[Porta]}
\]
A codificação reduz o tempo de atendimento em chamados de manutenção e facilita o rastreamento físico dos circuitos no laboratório.
